% ---- pg-364 （书页 352） ----

数据面作补充分析时, 常常必须确认它们近似是正态的.

最后几个主成分能帮助确认可疑的观测值. 各个观测值可以表达成 $\mathbf{S}$ 的全部特征向量 $\hat{\mathbf{e}}_1,\hat{\mathbf{e}}_2,\cdots,\hat{\mathbf{e}}_p$ 的线性组合
\[
\begin{aligned}
\mathbf{x}_j&=(\mathbf{x}_j'\hat{\mathbf{e}}_1)\hat{\mathbf{e}}_1+(\mathbf{x}_j'\hat{\mathbf{e}}_2)\hat{\mathbf{e}}_2+\cdots+(\mathbf{x}_j'\hat{\mathbf{e}}_p)\hat{\mathbf{e}}_p\\
&=\hat{y}_{j1}\hat{\mathbf{e}}_1+\hat{y}_{j2}\hat{\mathbf{e}}_2+\cdots+\hat{y}_{jp}\hat{\mathbf{e}}_p
\end{aligned}
\]
这样, 后几个主成分的大小就可以决定前几个主成分拟合观测值的好坏程度. 这就是说, $\hat{y}_{j1}\hat{\mathbf{e}}_1+\hat{y}_{j2}\hat{\mathbf{e}}_2+\cdots+\hat{y}_{j,q-1}\hat{\mathbf{e}}_{q-1}$ 与 $\mathbf{x}_j$ 的差为 $\hat{y}_{jq}\hat{\mathbf{e}}_q+\cdots+\hat{y}_{jp}\hat{\mathbf{e}}_p$, 它的长度平方是 $\hat{y}_{jq}^2+\cdots+\hat{y}_{jp}^2$. 可疑观测值常使坐标 $\hat{y}_{j1},\cdots,\hat{y}_{jp}$ 中至少有一个, 对这个平方长度的贡献较大(更一般的近似结果, 见补充 8A). 下面说法概括了这些观点.

\begin{enumerate}[label=\arabic*.,leftmargin=2em,itemsep=2pt]
\item 为了检查正态假定, 对前面几个主成分对, 构造散布图. 也可由各个主成分生成的样点值作出 $Q$-$Q$ 图.
\item 对后几个主成分构造散布图和 $Q$-$Q$ 图. 这些图可以帮助确认可疑的观测值.
\end{enumerate}

\noindent\textbf{例 8.7}\quad (对海龟数据作主成分图)

我们针对例 8.4 中讨论的雄海龟数据, 来解释如何作主成分图. 三个样本主成分是
\[
\begin{aligned}
\hat{y}_1&=0.683(x_1-4.725)+0.510(x_2-4.478)+0.523(x_3-3.703)\\
\hat{y}_2&=-0.159(x_1-4.725)-0.594(x_2-4.478)+0.788(x_3-3.703)\\
\hat{y}_3&=-0.713(x_1-4.725)+0.622(x_2-4.478)+0.324(x_3-3.703)
\end{aligned}
\]
其中 $x_1=\ln(\text{长度}), x_2=\ln(\text{宽度}), x_3=\ln(\text{高度})$.

\hyperref[fig:8.5]{图 8.5} 是 $\hat{y}_2$ 的 $Q$-$Q$ 图, \hyperref[fig:8.6]{图 8.6} 是 $(\hat{y}_1,\hat{y}_2)$ 的散布图. 对第一只海龟的观测值用圆圈圈出, 它位于散布图的右下方在 $Q$-$Q$ 图的右上方; 它可能是一个可疑点, 应该核对该点以查明是否有记录失误, 或者测查是否为结构异常. 删掉这只海龟数据后, 散布图接近椭圆. 其他的主成分对都不表明与正态情况有任何实质性的偏离. \hfill$\blacksquare$

\begin{mvfigure}
\begin{tikzpicture}
% ===== \hyperref[fig:8.5]{图 8.5}  Q-Q 图 =====
\begin{scope}[shift={(1.6,1.0)},x=1.0cm,y=26cm]
\draw[->] (-2.45,-0.05) -- (-2.45,0.078);
\draw[->] (-2.45,-0.05) -- (2.45,-0.05) node[below,font=\small] {$q_{(i)}$};
\node[left,font=\small] at (-2.5,0.0) {$\hat{y}_{(2)}$};
\foreach \y/\l in {-0.04/-0.04,0.0/0.0,0.04/0.04}{
  \draw (-2.52,\y) -- (-2.38,\y);
  \node[left,font=\scriptsize] at (-2.55,\y) {\l};}
\foreach \x/\l in {-2/-2,-1/-1,0/0,1/1,2/2}{
  \draw (\x,-0.052) -- (\x,-0.042);
  \node[below,font=\scriptsize] at (\x,-0.055) {\l};}
\foreach \x/\y in {-2.0/-0.033,-1.65/-0.032,-1.4/-0.030,-1.2/-0.027,-1.05/-0.025,-0.9/-0.022,-0.78/-0.020,-0.66/-0.018,-0.55/-0.015,-0.45/-0.012,-0.35/-0.010,-0.25/-0.008,-0.15/-0.006,-0.05/-0.004,0.05/-0.001,0.15/0.002,0.25/0.005,0.35/0.007,0.45/0.010,0.55/0.013,0.65/0.016,0.75/0.018,0.85/0.020,0.95/0.023,1.05/0.025,1.15/0.028,1.3/0.031,1.45/0.033}
  {\fill (\x,\y) circle (1.6pt);}
\node[draw,circle,inner sep=1.6pt] at (2.0,0.062) {};
\fill (2.0,0.062) circle (1.2pt);
\node[right,font=\small] at (2.05,0.062) {1};
\end{scope}
% ===== \hyperref[fig:8.6]{图 8.6}  散布图 =====
\begin{scope}[shift={(8.2,1.0)},x=46cm,y=7.6cm]
\draw[->] (-0.038,-0.30) -- (-0.038,0.30);
\draw[->] (-0.038,-0.30) -- (0.078,-0.30) node[below,font=\small] {$\hat{y}_2$};
\node[left,font=\small] at (-0.042,0.0) {$\hat{y}_1$};
\foreach \y/\l in {0.3/0.3,0.1/0.1,-0.1/-0.1,-0.3/-0.3}{
  \draw (-0.040,\y) -- (-0.032,\y);
  \node[left,font=\scriptsize] at (-0.042,\y) {\l};}
\foreach \x/\l in {-0.03/-0.03,-0.01/-0.01,0.01/0.01,0.03/0.03,0.05/0.05,0.07/0.07}{
  \draw (\x,-0.31) -- (\x,-0.29);
  \node[below,font=\scriptsize] at (\x,-0.315) {\l};}
\foreach \x/\y in {-0.030/0.06,-0.028/0.10,-0.024/0.02,-0.020/0.05,-0.017/0.02,-0.012/0.30,-0.011/0.02,-0.010/0.06,-0.006/0.06,-0.004/-0.005,-0.002/0.05,0.001/-0.15,0.003/-0.05,0.005/-0.12,0.008/-0.12,0.010/-0.12,0.012/0.18,0.014/0.19,0.016/0.14,0.020/0.10,0.022/0.06,0.024/0.22,0.026/0.08,0.028/-0.22,0.030/0.18,0.032/0.16,0.034/-0.14,0.036/-0.13,0.038/0.17}
  {\fill (\x,\y) circle (1.6pt);}
\node[draw,circle,inner sep=1.6pt] at (0.060,-0.290) {};
\fill (0.060,-0.290) circle (1.2pt);
\node[right,font=\small] at (0.062,-0.290) {1};
\end{scope}
\end{tikzpicture}
\phantomsection\label{fig:8.5}\mvcap{图 8.5\quad 雄海龟数据第二主成分 $\hat{y}_2$ 的 $Q$-$Q$ 图\qquad \phantomsection\label{fig:8.6}图 8.6\quad 雄海龟数据主成分 $\hat{y}_1$ 和 $\hat{y}_2$ 的散布图}
\end{mvfigure}

包含主成分的这一特征诊断, 可同样有效地运用到核查多元回归模型的假设. 事实上, 用任何估计法拟合完任意一个模型后, 对多元线性模型都要谨慎地考虑
\[
\text{残差向量}=(\text{观测值向量})-(\text{预测值(估计值)向量})
\]
或
